% Generated draft: requires mathematical and source review.
\documentclass[11pt]{article}
\usepackage[margin=25mm]{geometry}
\usepackage{fontspec}
\setmainfont{Latin Modern Roman}
\usepackage{amsmath,amssymb,mathtools,unicode-math}
\setmathfont{Latin Modern Math}
\usepackage{enumitem,needspace,xurl,hyperref}
\setlength{\parindent}{0pt}
\setlength{\parskip}{5pt}
\setlength{\emergencystretch}{3em}
\newcommand{\N}{\mathbb{N}}
\newcommand{\Z}{\mathbb{Z}}
\newcommand{\Q}{\mathbb{Q}}
\newcommand{\R}{\mathbb{R}}
\newcommand{\C}{\mathbb{C}}
\newcommand{\F}{\mathbb{F}}
\newcommand{\A}{\mathbb{A}}
\newcommand{\PP}{\mathbb{P}}
\newcommand{\E}{\mathbb{E}}
\newcommand{\eps}{\varepsilon}
\newcommand{\dd}{\,\mathrm{d}}
\begin{document}
\section*{UnsolvedMath ID: 700570}
\textbf{Source problem code:} catalogue entry
\subsection{Definitions and mathematical statement}
A topological group \(G\) is compactly generated if there exists a compact subset \(K\subseteq G\) such that \(G=\langle K\rangle\), where \(\langle K\rangle\) denotes the subgroup generated by \(K\).

The supplied problem asks whether every compactly generated, inductively prosoluble, locally compact group is prosoluble. Its logical form is
\[
\forall G\;\Bigl[\bigl(G\text{ is a locally compact topological group}\bigr)\land\bigl(\exists K\subseteq G\text{ compact}:G=\langle K\rangle\bigr)\land\mathsf{IndProSol}(G)\;\Longrightarrow\;\mathsf{ProSol}(G)\Bigr].
\]
Here \(\mathsf{IndProSol}(G)\) and \(\mathsf{ProSol}(G)\) are placeholders for the source's meanings of ``inductively prosoluble'' and ``prosoluble'', respectively. The supplied record does not define these terms or specify the convention concerning Hausdorffness for locally compact groups. Consequently, this is a conditional transcription of the intended implication; a fully specified mathematical statement requires those source conventions. In particular, no interpretation in terms of inverse limits or directed unions is imposed.
\subsection{Short English statement}
Must every compactly generated, inductively prosoluble, locally compact group be prosoluble? The meanings of the two solubility conditions require clarification from the source. The supplied registry labels the problem open; that label is not independently verified here.
\subsection{Sources}
\begin{itemize}
\item Local UnsolvedMath registry, permanent ID \(700570\), source code \texttt{catalogue entry}, title ``Kourovka Notebook Problem 9.37''. The supplied record attributes the problem to Yu. N. Mukhin and dates its proposal to 1984.
\item \emph{The Kourovka Notebook (Unsolved Problems in Group Theory)}, original problem 9.37, 9th issue (1984); supplied reference edition: 21st issue edition (2026), \texttt{arXiv:1401.0300v48}. The registry describes the collection as edited by E. I. Khukhro and V. D. Mazurov.
\item The supplied registry status is ``open''. The record supplies no dated literature-verification assessment; its bibliographic information and status have not been independently checked for this draft.
\end{itemize}
\end{document}
